\section*{Hoofdstuk \ref{ch:b2:living-soil} --- De levende bodem}

\begin{solution}{exo:b2:living-soil:1}
O: strooisel en gedeeltelijk vergaan organisch materiaal, waar de afbraak
begint. A: minerale \omterm{def:b2:living-soil:soil}{bodem} vermengd met \omterm{def:b2:living-soil:soil}{humus}, donker en met kruimelstructuur,
met de meeste wortels, organismen, voedingsstoffen en water. B: de horizont
van ophoping, waar klei, ijzeroxiden en carbonaten die van boven zijn
uitgespoeld worden afgezet. C: verweerd moedermateriaal, brokken gesteente in
een fijne matrix, waar nieuwe minerale \omterm{def:b2:living-soil:soil}{bodem} vrijkomt. R:
onverweerd vast gesteente.
\end{solution}

\begin{solution}{exo:b2:living-soil:2}
De negatieve ladingen aan het oppervlak van klei en \omterm{def:b2:living-soil:soil}{humus} houden uitwisselbare
kationen vast die wortels in ruil voor protonen kunnen opnemen. Zandkorrels
zijn kwarts, ongeladen en grof, met weinig oppervlak per gram
en vrijwel geen \omterm{def:b2:living-soil:soil}{humus}. Bekalken voegt calciumcarbonaat toe: het
calcium verdringt de waterstof- en aluminiumionen die de
uitwisselingsplaatsen van een zure \omterm{def:b2:living-soil:soil}{bodem} bezetten, het carbonaat neutraliseert
de zuurgraad, en de plaatsen worden opnieuw gevuld met een voedingskation.
\end{solution}

\begin{solution}{exo:b2:living-soil:3}
Afbrekers: bacteriën (\emph{Bacillus}) en schimmels
(\emph{Trichoderma}); microbiële grazers: \omterm{def:b2:unicellular-diversity:protist}{protisten} (amoeben) en
bacterie-etende nematoden; detritivoren: regenwormen, miljoenpoten,
springstaarten, pissebedden; predatoren: mijten, roofnematoden,
duizendpoten, loopkevers; en mollen of spitsmuizen aan de top. De
grazers, met een hogere $\mathrm{C:N}$ dan hun bacteriële prooi,
scheiden de overtollige stikstof uit als ammonium: ze mineraliseren de
stikstof die de bacteriën hadden vastgelegd.
\end{solution}

\begin{solution}{exo:b2:living-soil:4}
Arbusculair: de schimmel (een glomeromyceet) dringt de cellen van de
wortelschors binnen en vormt arbuscules; te vinden bij de meeste kruiden en
gewassen en bij tropische bomen. Ectomycorrhiza: de schimmel (basidiomyceten en
ascomyceten, veel ervan paddenstoelen) omhult de wortelpunt en
groeit tussen de cellen zonder ze binnen te dringen; te vinden bij bomen van
gematigde en boreale streken. In beide gevallen geeft de plant suiker en geeft
de schimmel fosfaat, stikstof, water, micronutriënten en enige
bescherming tegen ziekteverwekkers.
\end{solution}

\begin{solution}{exo:b2:living-soil:5}
$L^{*} = I/k = 4/0.8 = \qty{5}{t/ha}$; \omterm{def:b2:biogeochemical-cycles:reservoir}{verblijftijd} $1/k = 1.25$
jaar; \qty{90}{\%} verlies na $\ln 10/0.8 = 2.9$ jaar.
\end{solution}

\begin{solution}{exo:b2:living-soil:6}
Per \qty{100}{kg} gegeten koolstof houden de microben \qty{40}{kg} vast en
hebben ze $40/8 = \qty{5}{kg}$ stikstof nodig. Stro levert $100/80 =
\qty{1.25}{kg}$: immobilisatie van \qty{3.75}{kg} per \qty{100}{kg}
koolstof. Klaver levert $100/15 = \qty{6.7}{kg}$: mineralisatie
van \qty{1.7}{kg}.
\end{solution}

\begin{solution}{exo:b2:living-soil:7}
\qty{20}{t/ha} aarde per jaar door de wormen. De bovenste
\qty{20}{cm} weegt $0.2\times 10^{4}\times 1.3 = \qty{2600}{t/ha}$:
130 jaar voor één doorgang. De tien ton per acre van Darwin is
\qty{25}{t/ha}, dezelfde orde --- zijn ``om de paar jaar'' sloeg op de
oppervlakkige teellaag, enkele centimeters, niet op de hele bovengrond.
\end{solution}

\begin{solution}{exo:b2:living-soil:8}
Bacteriën: $10^{9}\times 10^{-12} = \qty{1}{mg}$ per gram. Schimmeldraden:
volume $\pi(2\times 10^{-4})^{2}\times 10^{4} = \qty{1.26e-3}{cm^3}$,
massa \qty{1.4}{mg} per gram. Per hectare ($2.6\times 10^{9}$ g):
\qty{2.6}{t} bacteriën en \qty{3.6}{t} schimmels.
\end{solution}

\begin{solution}{exo:b2:living-soil:9}
$H^{*} = 1.5/0.02 = \qty{75}{t/ha}$ koolstof. Geploegd: $1.5/0.04
= \qty{37.5}{t/ha}$; verlies \qty{37.5}{t/ha}, waarvan de helft verdwenen na
$\ln 2/0.04 = 17$ jaar.
\end{solution}

\begin{solution}{exo:b2:living-soil:10}
Fosfaat: $\sqrt{2\times 10^{-13}\times 8.64\times 10^{6}} =
\qty{1.3}{mm}$ in een seizoen --- een wortel op \qty{1}{cm} van de volgende
bereikt een achtste van de \omterm{def:b2:living-soil:soil}{bodem} ertussen, terwijl schimmeldraden op
\qty{0.1}{mm} van elkaar alles bereiken. Nitraat: \qty{4}{cm}, verder
dan de afstand tussen wortels: het komt vanzelf naar de wortel (en door
massastroming met het water), zodat mycorrhiza's weinig toevoegen voor nitraat
en veel voor fosfaat.
\end{solution}

\begin{solution}{exo:b2:living-soil:11}
Massa $0.25\times 10^{4}\times 1.3 = \qty{3250}{t/ha}$. Netto verlies
\qty{19}{t/ha} per jaar: levensduur 170 jaar. Verloren koolstof $20\times
0.03 = \qty{0.6}{t/ha}$ per jaar. Bij niet-kerende landbouw is het netto verlies
\qty{1}{t/ha}: 3250 jaar.
\end{solution}

\begin{solution}{exo:b2:living-soil:12}
Strooisel wordt in een jaar of twee vernieuwd, \omterm{def:b2:living-soil:soil}{humus} in vijftig, de minerale
A-horizont in duizenden (de vorming ervan met \qty{0.1}{mm} per jaar), en
zout hoopt zich in decennia op en wordt in decennia afgevoerd als er
drainage is. Een boer ziet het strooisel en het gewas binnen een seizoen
reageren en de \omterm{def:b2:living-soil:soil}{humus} binnen een loopbaan (een daling met een derde duurt
dertig jaar; een herstel even lang); de erosie van de minerale \omterm{def:b2:living-soil:soil}{bodem}
is in één mensenleven onzichtbaar en in tien onomkeerbaar; zout is
binnen een generatie zichtbaar en zonder drainage blijvend. De trage
variabelen van de \omterm{def:b2:living-soil:soil}{bodem} zijn die welke geen jaarrekening vastlegt.
\end{solution}

\begin{solution}{pb:b2:living-soil:1}
\textbf{1.} $0.25\times 10^{4}\times 1.3 = \qty{3250}{t/ha}$;
koolstof $\qty{81}{t/ha}$.
\textbf{2.} $81/12 = \qty{6.8}{t/ha}$ stikstof.
\textbf{3.} Bacteriën \qty{1}{mg/g}, $\qty{3.25}{t/ha}$; schimmeldraden
$\pi(2\times 10^{-4})^{2}\times 2\times 10^{4}\times 1.1 =
\qty{2.8}{mg/g}$, \qty{9}{t/ha}: microbiële biomassa ongeveer
\qty{12}{t/ha}.
\textbf{4.} Ongeveer \qty{14}{t/ha} levend materiaal (verse massa),
zo'n \qty{17}{\%} van de humuskoolstof --- als koolstof enkele procenten
daarvan, omdat organismen grotendeels uit water bestaan.
\textbf{5.} Microbiële koolstof ongeveer de helft van \qty{12}{t}, \qty{6}{t};
met twee vernieuwingen per jaar en \qty{40}{\%} retentie verbruiken ze
$2\times 6/0.4 = \qty{30}{t}$ koolstof en ademen ze \qty{18}{t/ha}
koolstof, \qty{66}{t} $\mathrm{CO_2}$ (een ruwe schatting: de
retentie en de vernieuwing zijn grof).
\textbf{6.} Drie keer de netto primaire productie van het gewas van
\qty{5}{t}: de schatting is te hoog, omdat de aannames (alle
biomassa twee keer vernieuwd, niets ervan in rust) royaal zijn;
de werkelijke \omterm{prop:b2:living-soil:humus}{bodemademhaling} onder een gewas ligt in de orde van de
eigen productie van het gewas, enkele ton koolstof per hectare per jaar. De
\omterm{def:b2:living-soil:soil}{bodem} ademt evenveel als de planten erboven.
\textbf{7.} Koolstof $5\times 0.45 = \qty{2.25}{t}$; stikstof $5\times
0.03 = \qty{150}{kg}$; $\mathrm{C:N} = 15$.
\textbf{8.} Het mineraliseert: de microben hebben $2250\times 0.4/8 =
\qty{112}{kg}$ stikstof nodig en het residu bevat er 150, zodat er netto
\qty{38}{kg/ha} vrijkomt.
\textbf{9.} $1 - e^{-0.5} = \qty{39}{\%}$ na 3 maanden; $1 - e^{-2}
= \qty{86}{\%}$ na een jaar.
\textbf{10.} $0.39\times 38 = \qty{15}{kg/ha}$.
\textbf{11.} Het gewas heeft \qty{150}{kg} nodig, grotendeels in zijn eerste
twee maanden; het residu geeft \qty{15}{kg} in drie maanden en
\qty{38}{kg} in totaal --- een aanvulling, geen voorziening. Zijn waarde ligt
ook in de stikstof die de klaver heeft gebonden en die in de eigen
\qty{150}{kg} van het residu zit, waarvan het meeste de \omterm{def:b2:living-soil:soil}{humus} ingaat en
over de jaren terugkomt.
\textbf{12.} Stro: koolstof \qty{2250}{kg}, stikstof $2250/80 =
\qty{28}{kg}$; de microben hebben \qty{112}{kg} nodig, zodat \qty{84}{kg/ha}
uit de \omterm{def:b2:living-soil:soil}{bodem} wordt geïmmobiliseerd --- meer dan de helft van de behoefte van
het gewas. De boer moet stikstof toevoegen met het stro, het ruim vóór het
zaaien onderploegen, of het aan het oppervlak laten liggen, waar het traag
afbreekt.
\textbf{13.} Geploegd: $1.2/0.025 = \qty{48}{t/ha}$; niet-kerend:
$1.2/0.015 = \qty{80}{t/ha}$.
\textbf{14.} $H(t) = 80 - 32\,e^{-0.015t}$: winst \qty{4.5}{t/ha}
na 10 jaar, \qty{17}{t/ha} na 50.
\textbf{15.} $\mathrm{d}H/\mathrm{d}t = 0.015\times 32 =
\qty{0.48}{t/ha}$ koolstof in het eerste jaar, \qty{1.8}{t}
$\mathrm{CO_2}$.
\textbf{16.} $0.48\times 1.5\times 10^{9} = \qty{0.7}{GtC}$ in het
eerste jaar, \qty{7}{\%} van de fossiele uitstoot; tegen het vijftigste jaar
is het tempo gedaald tot $0.48\,e^{-0.75} = \qty{0.23}{t/ha}$,
\qty{0.34}{GtC}.
\textbf{17.} De winst is de nadering van een nieuwe \omterm{def:b2:biogeochemical-cycles:reservoir}{stationaire toestand}: ze
neemt af tot nul wanneer $H$ \qty{80}{t/ha} bereikt, en de hele winst gaat
binnen enkele decennia terug naar de lucht als het veld opnieuw wordt geploegd
--- een bodemput is een eenmalige, omkeerbare overdracht, geen blijvende
compensatie van aanhoudende uitstoot.
\textbf{18.} $k_h = 0.02$: $H^{*} = \qty{60}{t/ha}$, een verlies van
\qty{20}{t/ha}, meer dan de winst van de eerste vijftig jaar niet-kerende landbouw: opwarming kan het
beheer tenietdoen.
\textbf{19.} Geploegd: $15 - 0.5 = \qty{14.5}{t/ha}$ per jaar; niet-kerend:
\qty{1}{t/ha}.
\textbf{20.} $3250/14.5 = 220$ jaar; $3250/1 = 3250$ jaar.
\textbf{21.} Koolstof $15\times 0.025 = \qty{0.38}{t/ha}$; stikstof
$0.38/12 = \qty{31}{kg/ha}$ per jaar.
\textbf{22.} De humusafbraak in de \omterm{def:b2:biogeochemical-cycles:reservoir}{stationaire toestand} onder de ploeg is $k_hH^{*} =
I = \qty{1.2}{t/ha}$ per jaar; erosie voegt daar nog een derde aan toe,
maar in tegenstelling tot de afbraak wordt ze niet door aanvoer gecompenseerd
--- ze put de voorraad uit.
\textbf{23.} Gedeeltelijk. Geërodeerde koolstof die in een sediment of een
stuwmeer wordt begraven, kan lang tegen afbraak beschermd zijn (een
begravingsput); maar de aggregaten breken tijdens het transport, veel van de
koolstof wordt onderweg geoxideerd, en de verloren vruchtbaarheid van het veld
wordt vervangen door kunstmest waarvan de productie koolstof uitstoot. Het
wereldwijde teken van het koolstofeffect van erosie staat nog ter discussie.
\textbf{24.} Strooisel $1/k = 0.5$ jaar; \omterm{def:b2:living-soil:soil}{humus} $1/k_h = 40$ jaar
onder de ploeg en 67 bij niet-kerende landbouw; minerale \omterm{def:b2:living-soil:soil}{bodem} $3250/0.5 =
6500$ jaar ten opzichte van de vorming, en 220 jaar ten opzichte van het
verlies onder de ploeg.
\textbf{25.} Koolstofvoorraad \qty{81}{t/ha} in de A-horizont; levensduur
220 jaar onder de ploeg, 3250 bij niet-kerende landbouw; het klaverresidu
levert ongeveer \qty{38}{kg/ha} minerale stikstof, waarvan 15 in de
eerste drie maanden.
\end{solution}
